\section{Evidência de IMP-IDEM-001 --- Etapa A (fundação canônica M7)}

\textbf{Feature:} \texttt{IMP-IDEM-001} --- receipt e retry universal de mutações
críticas.

\textbf{Estado anterior:} \texttt{DIVERGENTE}. Só \texttt{usuarios.ts} possuía
uma forma parcial de receipt, com hash \texttt{JSON.stringify(dados)}, sem
\texttt{tipo\_operacao}/\texttt{payload\_hash} canônicos, e \texttt{idOperacao}
opcional com geração de UUID novo.

\textbf{Estado depois:} \texttt{DIVERGENTE}. Esta rodada cria e comprova a
fundação canônica M7 reutilizável; nenhuma mutação de produto foi migrada e as
contagens da matriz não mudam.

\subsection{Contrato M7 derivado}

Fontes: Seção~7 (\texttt{sec:idempotencia-m7}), Seção~5
(\textit{Operações técnicas}) e o CUE \texttt{formal\_m7} (estados fechados,
\texttt{payload\_hash} hexadecimal de 64). Pontos confirmados:

\begin{itemize}
  \item identidade semântica \texttt{(uid, tipo\_operacao, payload\_hash)}; as
  três componentes são comparadas no retry;
  \item \texttt{hashPayload(tipoOperacao, payload) =
  SHA-256(tipoOperacao + "\textbackslash n" + canonicalize(payload))};
  \item \texttt{idOperacao} opaco, \texttt{[A-Za-z0-9\_-]\{1,128\}}, fora do
  hash; novo id só em nova intenção;
  \item \texttt{Operacoes/\{idOperacao\}}: \texttt{uid}, \texttt{tipo\_operacao},
  \texttt{payload\_hash}, \texttt{status} (PENDENTE/CONCLUIDA/FALHOU),
  \texttt{resultado}, \texttt{criado\_em}, \texttt{atualizado\_em};
  \item comando atômico grava \texttt{CONCLUIDA} na mesma transação; fluxo
  externo usa \texttt{PENDENTE}$\rightarrow$\texttt{CONCLUIDA}/\texttt{FALHOU};
  \item receipt corrompido falha fechado; retry não atualiza \texttt{timestamp}.
\end{itemize}

\subsection{Canonicalização}

A canonicalização ordena chaves de objeto por code unit recursivamente, ignora
\texttt{undefined} em objeto (equivalente a ausente), preserva a ordem de
arrays, distingue \texttt{null} de ausente e não aplica trim/caixa/arredondamento.
Tipos não suportados (não finitos, \texttt{bigint}, função, símbolo,
\texttt{undefined} isolado e objetos não planos como \texttt{Date}) falham
fechado, pois não pertencem ao domínio de payload JSON.

\subsection{Fundação implementada}

\begin{center}\small
\path{functions/src/idempotencia.ts}
\end{center}

\begin{itemize}
  \item \texttt{canonicalize(value)} --- impressão determinística do payload;
  \item \texttt{hashPayload(tipoOperacao, payload)} --- SHA-256 canônico global;
  \item \texttt{construirIdentidade(uid, tipoOperacao, payload)} --- monta a
  identidade semântica;
  \item \texttt{validarIdOperacao(idOperacao)} --- formato normativo;
  \item \texttt{resolverOperacaoTx(tx, idOperacao, identidade)} --- decide
  \texttt{NOVA}/\texttt{REPLAY}/\texttt{PENDENTE}/\texttt{FALHOU} ou rejeita
  reutilização incompatível com \texttt{already-exists}; receipt inconsistente
  falha fechado;
  \item \texttt{registrarOperacaoConcluidaTx}/\texttt{registrarOperacaoPendenteTx}
  --- escrita do receipt;
  \item \texttt{concluirOperacaoPendenteTx}/\texttt{falharOperacaoPendenteTx} ---
  transição \texttt{PENDENTE}$\rightarrow$\texttt{CONCLUIDA}/\texttt{FALHOU}.
\end{itemize}

\subsection{TDD: RED e GREEN}

RED contra o scaffold ingênuo (ordem de chave via \texttt{JSON.stringify},
comparação apenas de \texttt{uid}, sem validação de receipt):

\small
\begin{longtable}{@{}p{0.26\textwidth}p{0.55\textwidth}@{}}
\toprule
ID & Resultado \\
\midrule
\endhead
TEST-UNIT-M7-CANON-001 & FAIL (ordem de chave alterava o canonical) \\
TEST-UNIT-M7-CANON-002 & FAIL (objetos aninhados não ordenados) \\
TEST-UNIT-M7-CANON-009 & FAIL (tipos não suportados não falhavam fechado) \\
TEST-UNIT-M7-HASH-001 & FAIL (payload permutado gerava hash diferente) \\
TEST-INT-M7-RECEIPT-003 & FAIL (outro ator tratado como replay) \\
TEST-INT-M7-RECEIPT-004 & FAIL (outro tipo tratado como replay) \\
TEST-INT-M7-RECEIPT-005 & FAIL (outro payload tratado como replay) \\
TEST-INT-M7-RECEIPT-008 & FAIL (receipt corrompido não falhava fechado) \\
TEST-INT-M7-RECEIPT-009 & FAIL (\texttt{PENDENTE} tratado como nova) \\
TEST-INT-M7-RECEIPT-010 & FAIL (processador não implementado) \\
\bottomrule
\end{longtable}
\normalsize

Após a implementação, sob Node 20.19.6: unidade \texttt{TEST-UNIT-M7-*}
(18/18) e integração \texttt{TEST-INT-M7-RECEIPT-001--012} (12/12) aprovadas,
incluindo concorrência idêntica (\texttt{006}), intenções distintas
(\texttt{007}), \texttt{idOperacao} fora do hash (\texttt{011}) e
\texttt{FALHOU} terminal (\texttt{012}).

\subsection{Node 20 e regressão}

Evidência final sob Node \textbf{20.19.6} via
\path{nix shell github:nixos/nixpkgs/nixos-25.05#nodejs_20}
(\texttt{NODE20\_STRATEGY = TEMPORARY\_NIX\_ONLY}; Home Manager inalterado).

\small
\begin{longtable}{@{}p{0.52\textwidth}p{0.38\textwidth}@{}}
\toprule
Verificação (Node 20.19.6) & Resultado \\
\midrule
\endhead
\texttt{npm run build} & PASS \\
\texttt{npm run test:unit} & PASS (71/71, inclui 18 M7) \\
\texttt{TEST-INT-M7-RECEIPT-001--012} (emulator) & PASS (12/12) \\
\texttt{npm run test:integration} & 108 passam / 14 falham \\
\bottomrule
\end{longtable}
\normalsize

A baseline imediatamente anterior era 96/14; o delta são exatamente os 12 casos
M7 novos. As 14 falhas permanecem em \texttt{reagentes} (10),
\texttt{patrimonio} (2) e \texttt{roteiros} (2); nenhuma falha nova.

\subsection{Divergência registrada e limites}

\texttt{usuarios.ts} continua intocado e é evidência da divergência: usa
\texttt{Operacoes} com id derivado e \texttt{hash\_entrada =
JSON.stringify(dados)}, sem \texttt{tipo\_operacao}/\texttt{payload\_hash}
canônicos, e \texttt{idOperacao} opcional. Nenhum domínio de produto foi migrado;
nenhum processador/outbox externo foi implementado (apenas a taxonomia de
estados é suportada). A pendência externa de \texttt{convidarUsuario}
(Auth entre pré-checagem e commit), documentada na Etapa A.1, permanece
inalterada e é candidata natural à fundação M7/outbox posterior.

\textbf{Estado final:} \texttt{IMP-IDEM-001} continua \texttt{DIVERGENTE}. A
fundação M7 está estável e revisável isoladamente para receber a primeira
mutação real em rodada própria.
